\documentclass[11pt,a4paper]{article}
\usepackage[margin=23mm]{geometry}
\usepackage[T1]{fontenc}
\usepackage[utf8]{inputenc}
\usepackage{lmodern,microtype,amsmath,amssymb,booktabs,tabularx,array}
\usepackage{graphicx,xcolor,listings,enumitem}
\usepackage[section]{placeins}
\usepackage[colorlinks=true,linkcolor=blue!50!black,urlcolor=blue!50!black]{hyperref}
\usepackage{fancyhdr}
\pagestyle{fancy}
\fancyhf{}
\lhead{Statistical learning: worked notebook explanations}
\rfoot{\thepage}
\setlength{\headheight}{14pt}
\setlength{\parindent}{0pt}
\setlength{\parskip}{6pt}
\setlength{\emergencystretch}{2em}
\setlist{itemsep=3pt,topsep=4pt}
\lstset{language=Python,basicstyle=\small\ttfamily,breaklines=true,
 columns=fullflexible,keepspaces=true,showstringspaces=false,
 frame=single,rulecolor=\color{gray!35},backgroundcolor=\color{gray!4},
 commentstyle=\color{green!35!black},keywordstyle=\color{blue!65!black},
 aboveskip=8pt,belowskip=8pt}
\newcommand{\code}[1]{\texttt{\detokenize{#1}}}
\newcommand{\RSS}{\operatorname{RSS}}
\newcommand{\RMSE}{\operatorname{RMSE}}
\newcommand{\MAE}{\operatorname{MAE}}
\newcommand{\E}{\mathbb{E}}
\newcommand{\ind}[1]{\mathbf{1}\{#1\}}
\newcommand{\fig}[2]{\begin{figure}[htbp]\centering
\includegraphics[width=\linewidth,height=0.39\textheight,keepaspectratio]{figures/#1}
\caption{#2}\end{figure}}
\newcommand{\notebooktitle}[3]{\begin{center}{\LARGE\bfseries #1\par}\vspace{5pt}
{\large Problem statement and worked resolution\par}\vspace{5pt}
Companion to \code{#2}\par\end{center}
\textbf{How to use this explanation.} #3 Section numbers follow the notebook's numbered sections where present. Numerical values are rounded and refer to the recorded run, not universal performance claims.}
\newcommand{\methods}[1]{\subsection*{Main Python calls}#1}

\begin{document}
\notebooktitle{4A: Cluster geometry and stability}{04_A_cluster_geometry.ipynb}{This explanation derives the optimization steps, reconciles different quality measures, and shows what the software outputs do and do not establish.}
\section*{Problem statement and notation}
A clustering groups rows without a supplied response variable. Each $x_i\in\mathbb R^p$ is an observation;
$C_k$ is the set of row indices assigned to group k, and $\mu_k$ is its center.
K-means minimizes inertia, the within-cluster sum of squared distances,
\[W=\sum_{k=1}^K\sum_{i\in C_k}\|x_i-\mu_k\|^2.\]
This objective defines what the algorithm considers a good grouping. It does not define a universally correct grouping.
We use unscaled coordinates because both simulated features have the same units and specified noise variance.
There is no train/test accuracy: all rows are used to describe each dataset. Generating labels are reserved for a post-selection check.

\section{Why the Lloyd updates work}
For a fixed cluster C with mean $\bar x_C$, expand around that mean:
\[
\sum_{i\in C}\|x_i-\mu\|^2
=\sum_{i\in C}\|x_i-\bar x_C\|^2+|C|\|\mu-\bar x_C\|^2.
\]
The cross-term vanishes because centered residuals sum to zero. Thus the mean minimizes the loss for a fixed group.
For fixed centers, the loss separates over rows, so assigning each row to its nearest center minimizes its individual contribution.
These two facts prove that each alternating step cannot increase W. They do not prove a global minimum:
a different initial partition can converge to a different fixed point.

The notebook's array of squared distances has shape $(n,K)$: entry $(i,k)$ measures the cost of assigning row i to center k.
\code{argmin(axis=1)} chooses among centers for each row; using axis 0 would answer a different question.
An exact tie uses the first index. The teaching loop retains an old center for an empty group;
this is a declared convention, not a reproduction of every library recovery strategy.

For values $0,2,3,9,10,12$ and starting centers $(0,3)$:
\begin{center}\begin{tabular}{lrr}\toprule
Stage&Centers after update&W\\\midrule
Initial nearest assignments&$0,3$&167\\
First mean update&$0,7.2$&78.8\\
Second mean update&$5/3,31/3$&$28/3$\\
No further change&$5/3,31/3$&$28/3$\\\bottomrule\end{tabular}\end{center}
The first group initially contains only zero. After its partner center moves to 7.2, the boundary moves from 1.5 to 3.6,
so values 2 and 3 switch groups. The final midpoint is 6; assignments stop changing.
Each final group contributes $14/3$ to W. The repeated last value certifies stabilization under these updates, not global optimality in general.

\methods{\code{KMeans.fit} learns centers; \code{labels_} stores fitting-row assignments and \code{inertia_} stores their objective.
An explicit center array with \code{n_init=1} matches the starting condition. A small tolerance makes the numerical audit precise.
We compare centers, objective, and partition agreement; label names themselves need not match.}

\section{Read hierarchical merge heights in the correct units}
Agglomeration chooses the next pair of groups using a linkage criterion. For $\{0,1\}$ and $\{4,6\}$,
the cross-group distances are 4, 6, 3, and 5. Therefore the last single, complete, and average heights are
\[\min=3,\qquad\max=6,\qquad\text{mean}=18/4=4.5.\]
Ward's merge cost follows from the same centroid identity:
\[
\Delta(A,B)=W(A\cup B)-W(A)-W(B)
=\frac{|A||B|}{|A|+|B|}\|\bar x_A-\bar x_B\|^2.
\]
Here each group has size 2, their means are 0.5 and 5, and $\Delta=20.25$.
SciPy displays $\sqrt{2\Delta}=\sqrt{40.5}\approx6.363961$.
The first two Ward costs are 0.5 and 2, giving displayed heights 1 and 2.
Their costs sum to $22.75$, which is the squared error about the grand mean of all four rows.
This is an independent check on the interpretation of height.

\code{linkage} takes an observation matrix or a condensed distance vector. A square distance matrix passed as if it were
observations can silently cluster its rows as feature vectors. With Ward, use Euclidean observations as in the notebook.
\code{dendrogram} visualizes the merge table; moving a leaf horizontally does not change its merge history.

\section{Select K using an internal criterion}
The cloud simulation uses 600 rows from three equally probable Gaussian components, with noise standard deviation 0.6 per coordinate.
The candidate K values and the rule ``maximize silhouette, then smaller K on a tie'' are specified before inspecting generating labels.
For each candidate, 20 initializations seek a good inertia value. Results are:
\begin{center}\begin{tabular}{rrr}\toprule
K&Inertia&Mean silhouette\\\midrule
2&2717.382&0.573511\\3&441.976&0.776066\\4&388.267&0.595005\\
5&338.065&0.437117\\6&288.163&0.324521\\\bottomrule\end{tabular}\end{center}
The selected K is 3. Extra centers continue to reduce squared error but subdivide the well-separated clouds,
so silhouette decreases. This example makes the distinction between optimizing centers at a fixed K and selecting K visible.

For row i, silhouette compares its mean distance $a_i$ to other members of its group with the smallest mean distance $b_i$
to another group: $(b_i-a_i)/\max(a_i,b_i)$. The distance averages involve other observations, not just group centers.
Singletons have silhouette zero by the library convention. With fewer than two distinct groups, the usual mean silhouette is undefined.

\fig{04_A_cluster_geometry_2.png}{Inertia, silhouette, and the selected cloud partition answer different questions. The numerical labels are arbitrary; the black crosses are fitted centers.}
The generating-label adjusted Rand index (ARI) is 1 after selection. This confirms agreement with this simulation's components;
it did not participate in the choice of K. Selecting the highest ARI would turn the experiment into a label-informed task.

Twenty separate random single-start fits at K=3 have minimum and median inertia about 441.976 but maximum about 4153.276.
This substantial optimization failure explains why the restart count matters even on an apparently easy dataset.
\code{n_init=20} means 20 starts, not 20 Lloyd iterations; \code{max_iter} controls iterations within a start.

\section{Explain the ring counterexample}
Generate 600 angles and equally likely radii 1 and 3, then add coordinate noise of standard deviation 0.05.
Fix K=2 for this separate comparison. Results:
\begin{center}\begin{tabular}{lrr}\toprule
Ring method&Silhouette&Generating-label ARI\\\midrule
K-means&0.339157&$-0.001111$\\Single&0.153329&1.000000\\
Complete&0.343824&0.138083\\Average&0.319540&0.051850\\Ward&0.344317&0.113648\\\bottomrule\end{tabular}\end{center}
K-means creates nearest-center regions separated by linear boundaries: comparing squared distances to two centers cancels $x^Tx$.
Two such regions cannot represent an inner disk versus a surrounding annulus. The poor ARI is a mismatch of geometry, not a coding error.
Single linkage can follow short local connections around each ring and recovers the simulated groups here.
However, its lower silhouette is unsurprising: opposite points on one ring may be far apart while points on different rings are relatively close.

\fig{04_A_cluster_geometry_3.png}{All methods recover the separated clouds in this run. Their ring partitions differ because their criteria encode different group shapes.}
The silhouette ranking and label agreement ranking need not agree because they define quality differently.
Single linkage also has a weakness not measured by this clean run: a sparse bridge of points can create a chain joining two groups.
Its perfect ring ARI does not establish universal robustness.

\section{Interpret subsample stability correctly}
At selected K=3, draw ten independent 480-of-600 subsets without replacement and refit each time.
There are $\binom{10}{2}=45$ pairwise comparisons. Each compares only shared original row IDs;
overlap ranges from 377 to 393, with median 385. Every shared-row ARI is 1 in the recorded run.
This says the separated-cloud partition is insensitive to these row deletions at this fixed K.
It does not include uncertainty in K selection, because K is deliberately held fixed.

For two partitions, let $n_{ab}$ count rows in group a of one and group b of the other,
with marginal counts $n_{a\cdot}$ and $n_{\cdot b}$ and m shared rows. Define
\[
A=\sum_{a,b}\binom{n_{ab}}2,\quad B=\sum_a\binom{n_{a\cdot}}2,\quad
D=\sum_b\binom{n_{\cdot b}}2,\quad N=\binom m2.
\]
Then, outside degenerate cases handled by convention,
\[\operatorname{ARI}=\frac{A-BD/N}{(B+D)/2-BD/N}.\]
Counting pairs explains why renaming IDs has no effect, why chance adjustment can yield negative values,
and why comparing unrelated row orders is invalid. Pairwise stability scores reuse the same fits and are dependent;
the displayed quantiles are descriptive, not a confidence interval based on 45 independent experiments.

\section*{Documentation and practical checks}
The official \href{https://scikit-learn.org/1.8/modules/generated/sklearn.cluster.KMeans.html}{KMeans reference}
distinguishes \code{predict} (nearest-center IDs) from \code{transform} (distances to every center).
The \href{https://scikit-learn.org/1.8/modules/generated/sklearn.cluster.AgglomerativeClustering.html}{agglomerative reference}
provides fitting labels but no ordinary prediction method for new rows. Adding a centroid assignment rule to a hierarchy
would create an additional method, not recover an implicit library guarantee.
See also the \href{https://docs.scipy.org/doc/scipy/reference/generated/scipy.cluster.hierarchy.linkage.html}{linkage height conventions}
and the \href{https://scikit-learn.org/1.8/modules/generated/sklearn.metrics.adjusted_rand_score.html}{ARI reference}.

Keep seeds, feature units, distance, restart count, and K-selection rule in the report.
A stable result can still optimize an inappropriate geometry. Very high silhouette can be driven by a few isolated observations;
inspect group sizes and plots alongside the scalar score. The lesson is to connect the objective to the scientific or practical grouping question.
\end{document}
